% NSF PESOSE proposal (NSF 26-506) -- Project Description, TRACK 3 (Improving safety, security, privacy).
% Track 3: max $1,500,000 over up to 2 years. Project Description limit: 15 pages.
% Title prefix MUST be "PESOSE: Track 3: <Title>". Last line of the Project Summary MUST be a
% prioritized "Keywords:" list (2-5, semicolon-separated). No URLs in the Project Description body --
% point to the open-source product via an in-line citation + a References Cited entry.
%
% Track 3 Project Description MUST address (see solicitation):
%   - Societal, National, and Economic Impact (+ what depends on the OSE)   -> impact.tex
%   - Targeted Vulnerabilities and Risks (technical + socio-technical)      -> vulns.tex
%   - Development Plan (with yearly milestones; memory-safe langs, SBOMs)   -> devplan.tex
%   - Evaluation Plan (metrics/benchmarks, ideally validated with users)    -> eva.tex
% Plus every PESOSE proposal: pointer to the existing public product, its current dev/testing model,
% dissemination, user base, contributor base; the problem; strong team-qualification justification.
% Page budget (15 pages):
%   Overview/Objectives ~2 | The FreeRTOS OSE (product status) ~1.5 | Impact ~1.5
%   Vulnerabilities & Risks ~3 | Development Plan ~5 | Evaluation ~1 | Broader Impacts/Timeline/Prior ~1

\documentclass[11pt,oneside,onecolumn,letterpaper]{article}
%\usepackage[OT1]{fontenc}
\usepackage{times}
\usepackage[margin=1in]{geometry}
%\usepackage{amsmath,amssymb,amsfonts}
\usepackage{booktabs}
\usepackage{graphicx}
\usepackage{latexsym}
%\usepackage{subfigure}
\usepackage{wrapfig}
\usepackage{amsthm}
\usepackage{amsmath,amssymb,amsfonts}
\usepackage[hyphens]{url}
\usepackage{pifont}
\usepackage{color}
\usepackage{colortbl}
\usepackage[dvipsnames]{xcolor}
\usepackage[ruled,vlined]{algorithm2e}
\usepackage[square, comma, sort&compress, numbers]{natbib}
\usepackage{pgfgantt}
\usepackage{enumitem}
\usepackage{subcaption}
\usepackage{titlesec}
\usepackage{xspace}
\usepackage{arydshln}
\usepackage{circledsteps}
\usepackage{listings}
\usepackage{stfloats}
\usepackage{comment}
\usepackage{fmtcount}
\usepackage{lipsum}
\usepackage{soul}
\usepackage[textsize=tiny]{todonotes}
\usepackage{multirow}
\usepackage{multicol}
\usepackage{threeparttable}
\usepackage{tcolorbox}
\usepackage{minted}
\usepackage{graphicx}
\usepackage{tikz}
\usetikzlibrary{positioning,arrows.meta,calc,fit}

\newcommand{\rulebox}[1]{
	\begin{tcolorbox}[
		%width=4in,
		boxsep=0pt,
		left=2pt,
		right=2pt,
		top=1pt,
		bottom=1pt,
		boxrule=0.5pt,
		]
		#1
	\end{tcolorbox}
}

\newcommand*\circled[1]{\tikz[baseline=(char.base)]{
		\node[shape=circle,draw,inner sep=0.7pt] (char) {#1};}}

\titlespacing\section{0pt}{4pt plus 2pt minus 2pt}{6pt plus 2pt minus 2pt}
\titlespacing\subsection{0pt}{4pt plus 2pt minus 2pt}{2pt plus 2pt minus 2pt}
\titlespacing\subsubsection{0pt}{4pt plus 4pt minus 2pt}{2pt plus 2pt minus 2pt}

\titleformat{\section}
{\fontsize{12}{15}\bfseries}{\thesection}{1em}{}

\titleformat{\subsection}
{\fontsize{11}{15}\bfseries}{\thesubsection}{1em}{}

% Inline review/comment macro (highlighted). Wrap notes you do NOT want in the final PDF.
\newcommand{\ziming}[1]{%
  \begingroup
  \definecolor{hlcolor}{RGB}{20, 255, 20}\sethlcolor{hlcolor}%
  \textcolor{black}{\hl{\textit{\textbf{Ziming:} #1}}}%
  \endgroup
}

% Salient placeholder/TODO marker -- bright red bold so unfilled gaps are impossible to miss
% in the compiled PDF. Remove before submission.
\definecolor{todored}{rgb}{0.85,0,0}
\newcommand{\TODO}[1]{\textcolor{todored}{\textbf{[TODO: #1]}}}

% Placeholder for a not-yet-available photo/figure: a red-framed gray box with a label.
\newcommand{\photoplaceholder}[2][0.3\textwidth]{%
  \fcolorbox{red}{gray!12}{\parbox[c][2.6cm][c]{#1}{\centering\footnotesize\textcolor{red}{\textbf{[PHOTO: #2]}}}}%
}

% ---------------------------------------------------------------------------
% TODO(pesose): name the effort/managing-effort this proposal introduces.
% \sysname is a working placeholder name -- rename to the final project name.
% ---------------------------------------------------------------------------
\newcommand{\sysname}{{\scshape Sentry}\xspace}      % TODO(pesose): pick the final effort name
\newcommand{\ose}{\textsc{FreeRTOS}\xspace}           % the open-source product / ecosystem being secured

% Development-Plan activity titles. Arc: discover -> prevent/harden -> detect & recover.
% This embeds the team's core "find the bugs" goal (Activity 1) inside a full Track-3 hardening agenda,
% as the solicitation requires (prevention, detection, recovery). Refine wording as the proposal matures.
\newcommand{\actone}{Continuous, Scalable Vulnerability Discovery for the \ose Codebase}
\newcommand{\acttwo}{Prevention and Hardening: Memory Safety, Secure-Development Practices, and SBOMs}
\newcommand{\actthree}{Detection, Response, and Recovery for the \ose Ecosystem}

\newcommand{\cmark}{\ding{51}}%
\newcommand{\xmark}{\ding{55}}%

\definecolor{codegreen}{rgb}{0,0.6,0}
\definecolor{codegray}{rgb}{0.5,0.5,0.5}
\definecolor{codepurple}{rgb}{0.58,0,0.82}
\definecolor{backcolour}{rgb}{0.95,0.95,0.92}

\pgfkeys{/csteps/inner xsep=1.5pt}
\pgfkeys{/csteps/inner ysep=1.5pt}

\colorlet{punct}{red!60!black}
\definecolor{background}{HTML}{EEEEEE}
\definecolor{delim}{RGB}{20,105,176}
\colorlet{numb}{magenta!60!black}

\lstdefinestyle{mystyle}{
    backgroundcolor=\color{backcolour},
    commentstyle=\color{codegreen},
    keywordstyle=\color{magenta},
    numberstyle=\tiny\color{codegray},
    stringstyle=\color{codepurple},
    basicstyle=\ttfamily\tiny,
    breakatwhitespace=false,
    breaklines=true,
    captionpos=b,
    keepspaces=true,
    numbers=none,
    numbersep=5pt,
    showspaces=false,
    showstringspaces=false,
    showtabs=false,
    tabsize=2,
}

\newcounter{goal}
\newcommand{\goal}[1]{
  \stepcounter{goal}
  \textbf{G\padzeroes[1]\decimal{goal}. #1}
}

\lstset{style=mystyle}

\graphicspath{{./figs/}}

\begin{document}

\pagenumbering{gobble}

\begin{center}
{\large{\bf
% Title MUST start with "PESOSE: Track 3:" per the solicitation.
PESOSE: Track 3: \sysname: Securing the \ose{} Open-Source Ecosystem through Continuous Vulnerability Discovery, Hardening, and Recovery}}
\end{center}

\section{Overview and Objectives}

% TODO(pesose): Motivate the problem in 2-3 paragraphs.
%  - FreeRTOS is one of the most widely deployed real-time operating systems for microcontrollers,
%    shipping inside billions of embedded/IoT devices across critical infrastructure, medical,
%    industrial control, automotive, and consumer products. It is MIT-licensed and community-driven.
%  - A vulnerability in FreeRTOS (or its libraries / supply chain) propagates to every dependent
%    device, many of which are unpatchable in the field -> outsized societal/national risk.
%  - The problem this Track 3 effort addresses: systematically discover safety/security/privacy
%    vulnerabilities in the FreeRTOS ecosystem and reduce them through prevention, detection, and
%    recovery -- with fixes validated against real deployments and upstreamed.
%  - State the overarching goal of this proposal in one bold sentence.

% NOTE (solicitation rule): NO URLs in the Project Description body. Point to the open-source
% product with an in-line citation that resolves to a References Cited entry, e.g.:
The open-source product at the center of this proposed ecosystem effort is the \ose{} real-time operating system kernel and its associated libraries~\cite{freertos}.

\subsection{Proposed Effort}

% TODO(pesose): one paragraph overview + an overview figure (understand the ecosystem -> discover
% vulnerabilities -> prevent/harden -> detect & recover -> upstream). Reuse the tikz overview style
% from the lab's template if helpful.
This proposal pursues a three-activity agenda, \sysname, that follows a \emph{discover $\rightarrow$ prevent/harden $\rightarrow$ detect \& recover} arc, all fed back upstream into the \ose{} ecosystem.

\noindent\textbf{Activity 1: \actone.}
% The team's core capability: find the bugs. Continuous, scalable discovery (fuzzing, static/dynamic
% analysis, symbolic execution) across the kernel and libraries. This is the engine, not the deliverable.

\noindent\textbf{Activity 2: \acttwo.}
% Turn findings into durable prevention: memory-safety hardening, secure-development practices,
% software bills of materials (SBOMs), and supply-chain integrity for the ecosystem.

\noindent\textbf{Activity 3: \actthree.}
% Detection, response, and recovery: CI security gates, monitoring, coordinated disclosure, and a
% sustainable upstreaming/recovery process with the FreeRTOS maintainer community.

\subsection{Relevance to PESOSE (Track 3)}

% TODO(pesose): tie explicitly to the solicitation. PESOSE Track 3 "addresses vulnerabilities in
% open-source products and infrastructure" -- technical AND socio-technical (supply chain, insider,
% poisoned ML). Stress that this is ecosystem HARDENING (prevention/detection/recovery), validated
% with users and upstreamed -- not standalone bug-finding research.
This effort directly serves PESOSE Track 3, which strengthens the security, privacy, and resilience of critical open-source ecosystems by identifying vulnerabilities, mitigating risks, and building sustainable prevention, detection, and recovery processes.

\subsection{Team Qualifications}

% TODO(pesose): make the case that the team is qualified (REQUIRED by the solicitation), AND that
% it can get fixes ADOPTED upstream (the key Track-3 risk for an outside team). If there is any
% relationship with the FreeRTOS maintainer/AWS community, prior accepted PRs, or an MoU, state it.
The team has a sustained track record in embedded and microcontroller (MCU) systems security~\cite{tan24sok, nino24firware, ma2023dac, zhao24tee}, precisely the layer on which \ose{} runs.
\textbf{PI Zhao}, an Associate Professor at the Khoury College of Computer Sciences at Northeastern University, leads a group that systematically studies MCU/IoT firmware security, including crawling and analyzing thousands of real-world MCU firmware samples to identify vulnerabilities~\cite{nino24firware}.
The team also brings deep hands-on offensive and defensive embedded experience: PI Zhao has led teams in the MITRE Embedded Capture-the-Flag (eCTF) competition for six consecutive years~\TODO{confirm years/placements and cite}, giving an attacker's-eye understanding of exactly the vulnerability classes this effort targets.
\textbf{Co-PI Engin Kirda}, a Professor at Northeastern, brings extensive expertise in systems security, malware analysis, and large-scale program analysis~\TODO{cite Kirda's most relevant works and add 1-2 sentences on his role (discovery/analysis tooling for Activity~1)}.
% TODO(pesose): state the division of labor between PI Zhao and co-PI Kirda explicitly.
The effort is further supported by a validation collaborator, \textbf{Prof.\ Aanjhan Ranganathan} (Northeastern), who operates a drone/IoT security testbed built on \ose-based platforms (e.g., the Crazyflie nano-quadcopter and ESP32-class nodes), providing real-deployment validation for discovered-and-hardened vulnerabilities and a letter of collaboration as a \ose{} user.
% TODO(pesose): confirm Ranganathan's exact FreeRTOS-based platforms; he is a COLLABORATOR (letter +
% validation), not a co-PI. This also helps satisfy the 3-5 letters-of-collaboration requirement and
% the Track-3 "validation with existing users" evaluation requirement.


\section{The \ose{} Open-Source Ecosystem}

% REQUIRED by the solicitation for EVERY PESOSE proposal:
%   "details on the current status of the open-source product development and testing model,
%    methods of dissemination, user base, and contributor base."
% This section supplies exactly that for FreeRTOS. Keep all URLs out -> cite \cite{freertos} etc.

\noindent\textbf{The product.}
% TODO(pesose): what FreeRTOS is -- a real-time kernel + libraries (FreeRTOS+TCP, FreeRTOS+FAT,
% coreMQTT, corePKCS11, etc.); MIT license; portable across many MCU architectures.
\ose{} is an MIT-licensed real-time operating system kernel for microcontrollers and small microprocessors, together with a set of companion libraries~\cite{freertos}.

\noindent\textbf{Development and testing model.}
% TODO(pesose): how FreeRTOS is developed and tested today -- GitHub-based, AWS stewardship,
% maintainers, release cadence, existing CI, existing test/coverage practices, MISRA conformance,
% existing static-analysis usage. Be accurate; cite sources in References Cited.

\noindent\textbf{Methods of dissemination.}
% TODO(pesose): GitHub repositories, official site, package ecosystems, vendor SDKs that bundle
% FreeRTOS (e.g., chip-vendor toolchains), Amazon-FreeRTOS / FreeRTOS LTS distributions.

\noindent\textbf{User base.}
% TODO(pesose): scale and breadth of deployment (billions of devices; sectors: industrial control,
% medical, automotive, consumer IoT, aerospace). Quantify with citable figures where possible.

\noindent\textbf{Contributor base.}
% TODO(pesose): who develops FreeRTOS -- AWS team + external contributors, GitHub contributor counts,
% governance (note: this informs the upstream-adoption strategy in Activity 3).

\noindent\textbf{The problem being addressed.}
% TODO(pesose): crisp statement of the safety/security/privacy problem in the FreeRTOS ecosystem
% that this effort will reduce (leads into the Vulnerabilities & Risks section).


\section{Societal, National, and Economic Impact}

% Track 3 REQUIREMENT (solicitation): "Explain the broader impacts of the OSE. Describe any other
% products or systems that depend on the secure and privacy-preserving function of the OSE."

\noindent\textbf{What depends on \ose{}.}
% TODO(pesose): map the dependency surface -- the products, SDKs, and critical-infrastructure systems
% that embed FreeRTOS (industrial control, medical devices, automotive ECUs, smart-grid endpoints,
% aerospace, consumer IoT). Emphasize that many are long-lived and hard/impossible to patch in the
% field, so an upstream vulnerability has outsized, durable blast radius.

\noindent\textbf{National and economic significance.}
% TODO(pesose): tie to national importance (critical infrastructure, supply-chain security executive
% guidance) and to the economic value of open-source (the solicitation cites >$13T/yr for OSS overall).
% A single hardened, trusted RTOS ecosystem de-risks an enormous installed base.

\noindent\textbf{Why this is not adequately addressed today.}
% TODO(pesose): make the Track-3 review-criterion case that this is "an issue of significant societal
% or national importance" that is "not currently being adequately addressed" -- current FreeRTOS
% security relies on limited maintainer bandwidth; systematic, continuous, ecosystem-wide hardening
% with provable coverage is missing.


\section{Targeted Vulnerabilities and Risks}
\label{sec:vulns}

% Track 3 REQUIREMENT (solicitation): "Specify the classes of vulnerabilities to be addressed and
% the expected impact of mitigating them. Discuss attack methods being targeted, including technical
% risks and socio-technical risks. Include any known prior incidents or patterns of exploitation."
% Also satisfies the Track-3 review criterion: "clearly describe the vulnerability landscape."

\ose{} is written in C and runs with no memory-management unit and few hardware protections on the typical microcontroller, so a single memory-safety defect in widely reused code can yield remote code execution on millions of fielded, hard-to-patch devices.
We therefore prioritize the vulnerability classes by \emph{blast radius}: the network-facing, attacker-reachable parsing code shared across the ecosystem, and the components whose compromise propagates to the most dependents.
Three components are \emph{in scope} for this proposal: (i) the \ose+TCP/IP networking stack, the historical CVE hotspot and the most directly attacker-reachable code; (ii) the connectivity libraries layered on top of it (coreMQTT, coreHTTP, and the coreJSON parser they rely on), which parse untrusted protocol data; and (iii) the kernel concurrency primitives (the scheduler, queues, stream/message buffers, and ISR/critical-section handling), whose real-time-specific defects are poorly served by existing tooling.
% TODO(pesose): justify the prioritization in 1-2 sentences. Cryptographic/key-management and OTA
% components are adjacent and acknowledged below, but are NOT primary targets of this effort.

\subsection{Technical Vulnerability Classes}

\noindent\textbf{Memory-safety defects in attacker-reachable C code.}
The highest-risk surface is the network-facing stack, which parses untrusted input directly on the device.
This is not hypothetical: in 2018, Zimperium zLabs disclosed 13 vulnerabilities in the FreeRTOS+TCP/IP stack used by FreeRTOS and AWS FreeRTOS, including remote code execution, information leak, and denial of service reachable by a network adversary~\cite{zimperium2018freertos}~\TODO{confirm the exact CVE IDs (CVE-2018-16522..16528, 16599..16603) and cite each}.
Subsequent out-of-bounds and overflow CVEs in the same stack's DNS, DHCP, ARP, and ND parsing show this is a \emph{recurring} class, not a one-time event~\TODO{cite the later FreeRTOS+TCP CVEs}.
Mitigating this class removes the most directly weaponizable bugs in the ecosystem.

\noindent\textbf{Integer and length-handling errors.}
% TODO(pesose): arithmetic/length bugs (e.g., in packet/length fields) that feed the overflows above.

\noindent\textbf{Concurrency and RTOS-specific defects (primary target).}
A real-time kernel exposes defect classes that conventional memory-safety tooling largely ignores: data races and unsynchronized access to queues, stream buffers, and message buffers; priority inversion; and incorrect use of critical sections and interrupt-service routines.
Because these depend on scheduling and interrupt timing rather than on a single input, they are hard to find with input fuzzing alone, making them both high-impact and an opportunity for a distinctive RTOS-aware analysis contribution.
% TODO(pesose): name the concrete kernel objects/APIs in scope and the RTOS-aware analysis approach
% (tie to Activity 1).

\noindent\textbf{Cryptographic and OTA components (acknowledged, not primary).}
% TODO(pesose): briefly acknowledge corePKCS11 / TLS integration / OTA signature verification as
% adjacent confidentiality-integrity surface, but state they are out of this proposal's primary scope.

\subsection{Socio-Technical Risks}

\noindent\textbf{Supply-chain and fragmentation risk.}
Chip vendors routinely fork and bundle \ose{} inside their SDKs, where it drifts from upstream; without robust software bills of materials (SBOMs) and provenance, a downstream integrator cannot tell whether a given upstream CVE affects their shipped firmware.
% TODO(pesose): expand -- dependency provenance, build/release integrity, the SBOM gap.

\noindent\textbf{Maintainer bandwidth and bus-factor risk.}
% TODO(pesose): the human/governance risk -- limited maintainer capacity relative to the size of the
% dependent base; how this effort adds sustainable security capacity rather than one-off patches.

\noindent\textbf{Untrusted contributions and ports.}
% TODO(pesose): risk from third-party ports/contributions; (if relevant) poisoned dependencies.

\subsection{Attack Methods and Prior Incidents}

\noindent\textbf{Remote network exploitation.}
The 2018 FreeRTOS+TCP cluster~\cite{zimperium2018freertos} is the archetypal attack method: a remote, unauthenticated adversary sends crafted packets to reach the parsing bugs above.

\noindent\textbf{Physical fault injection (capability noted; out of scope).}
The microcontrollers \ose{} runs on are also exposed to physical fault injection (e.g., voltage/clock glitching), and the PIs bring hands-on fault-injection experience from six years of the MITRE eCTF competition.
We note this capability for completeness and to establish the team's depth at the hardware/software boundary, but physical fault injection is \emph{explicitly out of scope} for this proposal, which focuses on discovering and reducing software vulnerabilities in the \ose{} codebase and its libraries.
% TODO(pesose): keep this framing tight -- one or two sentences; do not let it expand the scope.


\section{Development Plan}

% Track 3 REQUIREMENT (solicitation): "Provide a detailed plan for addressing vulnerabilities,
% including key milestones for each year of the award. For software-focused OSEs, note technical
% considerations such as using memory-safe languages or software bills of materials."
% Structure: three activities = discover -> prevent/harden -> detect & recover, each with yearly
% milestones. (Year-by-year milestones are also listed near the Timeline in ProjectDescription.tex.)

\subsection{Activity 1: \actone}
\label{sec:act1}

\noindent\textbf{Objective.}
% TODO(pesose): continuously and scalably surface exploitable defects in the FreeRTOS kernel and
% libraries -- the team's core "find the bugs" capability, framed defensively.

\noindent\textbf{Approach.}
% Scope (locked): FreeRTOS+TCP, the coreMQTT/coreHTTP/coreJSON connectivity parsers, and the kernel
% concurrency primitives. Crypto/OTA and physical fault injection are out of primary scope (see vulns.tex).
% TODO(pesose): the discovery toolbox and how it is applied to these targets:
%  - Coverage-guided fuzzing of the parsing-heavy targets (FreeRTOS+TCP, coreMQTT, coreHTTP, coreJSON),
%    with harnesses + emulation (e.g., QEMU/renode) so campaigns run at scale off-device.
%  - Static analysis / sanitizers / symbolic execution for memory-safety and length/integer bugs.
%  - RTOS-aware analysis for scheduler/queue/stream-buffer/ISR/critical-section races -- the
%    distinctive, under-tooled class.
%  - Triage + reproducible PoCs, validated on real hardware where needed (team testbed).

\noindent\textbf{Year 1 / Year 2 milestones.} \TODO{e.g., Y1: discovery infra + harnesses for N libraries, first confirmed-and-reported vulnerabilities; Y2: continuous campaign + measured coverage targets.}

\subsection{Activity 2: \acttwo}
\label{sec:act2}

\noindent\textbf{Objective.}
% TODO(pesose): turn findings into durable PREVENTION across the ecosystem.

\noindent\textbf{Approach.}
% TODO(pesose): solicitation explicitly mentions memory-safe languages and SBOMs:
%  - Selective memory-safety hardening of hot/critical components (e.g., Rust components or
%    safe-C transformations / bounds checking) where it can be upstreamed without breaking the
%    MCU footprint -- discuss feasibility honestly given FreeRTOS's resource constraints.
%  - Secure-development practices: CISA/NSA secure-software-supply-chain guidance and the OpenSSF
%    best practices (the solicitation explicitly encourages these). Cite in References Cited.
%  - Software Bills of Materials (SBOMs) and provenance for FreeRTOS releases and vendor bundles.

\noindent\textbf{Year 1 / Year 2 milestones.} \TODO{e.g., Y1: SBOM + secure-dev baseline, hardening prototype for one component; Y2: upstreamed hardening for K components.}

\subsection{Activity 3: \actthree}
\label{sec:act3}

\noindent\textbf{Objective.}
% TODO(pesose): detection, response, and recovery -- the resilience half of Track 3.

\noindent\textbf{Approach.}
% TODO(pesose):
%  - Detection: security gates in CI (fuzzing/sanitizers/regression oracles) that run continuously.
%  - Response: coordinated disclosure + patch pipeline with the FreeRTOS maintainer community.
%  - Recovery: a sustainable, documented process so the ecosystem keeps benefiting after the award
%    (downstream notification, LTS backports, SBOM-driven impact analysis for dependents).

\noindent\textbf{Year 1 / Year 2 milestones.} \TODO{e.g., Y1: CI security gate piloted; Y2: end-to-end disclosure-to-recovery process demonstrated and handed to maintainers.}

\subsection{Upstream Adoption Strategy}
% TODO(pesose): CRITICAL for an outside team. How findings/hardening actually LAND in FreeRTOS:
% relationship with maintainers/AWS, contribution workflow, prior accepted PRs, letters of
% collaboration from FreeRTOS users/contributors (3-5 required as a separate supplementary doc).


\section{Evaluation Plan}

% Track 3 REQUIREMENT (solicitation): "Describe how success will be measured. Include metrics, tools,
% and benchmarks for assessing progress and effectiveness. Ideally, the evaluation plan should
% incorporate testing and validation opportunities with existing users."
% Also a Track-3 review criterion: "a specific, actionable list of milestones and evaluation plan."

\noindent\textbf{Activity 1 (Discover).}
% TODO(pesose): metrics -- number/severity of confirmed vulnerabilities; code/edge coverage achieved
% on target libraries; reproduction rate; time-to-find vs. baselines. Benchmarks: known-bug corpora /
% injected-bug benchmarks for FreeRTOS libraries. \TODO{set concrete targets and name baselines}

\noindent\textbf{Activity 2 (Prevent/Harden).}
% TODO(pesose): metrics -- vulnerability classes eliminated by construction; hardening overhead
% (code size, RAM, run-time) on representative MCU targets (must stay within embedded budgets);
% SBOM coverage of the release; secure-dev practice conformance. \TODO{set concrete targets}

\noindent\textbf{Activity 3 (Detect/Recover).}
% TODO(pesose): metrics -- CI gate catch-rate on regression/injected bugs; time-from-disclosure-to-
% fix; fraction of fixes upstreamed and adopted; downstream-notification coverage via SBOM.

\noindent\textbf{Validation with existing users.}
A primary validation channel is the \ose-based drone/IoT testbed operated by collaborator Prof.\ Ranganathan (e.g., Crazyflie and ESP32-class platforms), on which discovered-and-hardened vulnerabilities are reproduced and the hardening's footprint/overhead is measured on real deployments.
% TODO(pesose): detail the validation protocol on Ranganathan's testbed; add any other FreeRTOS
% users/contributors who will validate (ties to the 3-5 required letters of collaboration).
\TODO{name additional validation partners / real deployments}


\section{Broader Impacts}
\label{sec:bi}

% PESOSE is reviewed on the two NSB criteria (Intellectual Merit + Broader Impacts) PLUS the
% track-specific criteria. Keep a genuine Broader Impacts section. The ecosystem-security mission
% itself is a large broader impact; below are the education/workforce/outreach components the team
% reliably delivers. Also map to the PESOSE "Measures of Success": Discovery & Innovation
% (new tools/techniques, datasets), STEM Education & Workforce (participants into STEM), Research
% Infrastructure (new infrastructure built).

\textbf{Securing critical infrastructure that runs on \ose{}.}
% TODO(pesose): the primary broader impact -- hardening a RTOS embedded in medical, industrial,
% automotive, and IoT systems protects public safety and the economy at ecosystem scale.

\textbf{Open Tools, Benchmarks, and Datasets (Measure of Success: Discovery \& Innovation).}
% TODO(pesose): the discovery harnesses, SBOMs, vulnerability/benchmark datasets, and CI security
% gates produced here will themselves be released open-source for the FreeRTOS community and beyond.

\textbf{Workforce, CTF, and Education (Measure of Success: STEM Workforce).}
The PIs have a strong embedded-security education and competition record that this effort extends.
% TODO(pesose): MCU CTF platform, courses (Software Vulnerabilities and Security; IoT Security;
% Trusted Computing), CactiLab YouTube \cite{youtubelist} \TODO{update subscriber/view stats}, and
% the MITRE eCTF team "Cacti" \TODO{years/placements}. Tie student training to FreeRTOS hardening.

\textbf{Broadening Participation.}
% TODO(pesose): undergraduate, high-school, and underrepresented-student involvement (the team's
% prior record, e.g., high-school advisees and the 300+ -student Cacti team).


\section{Project Timeline and Milestones}

\begin{wrapfigure}{r}{0.4\textwidth}
	\vspace{-1cm}
	\definecolor{barblue}{RGB}{153,204,254}
	\definecolor{groupblue}{RGB}{51,102,254}
	\definecolor{linkred}{RGB}{165,0,33}
	\centering
	\begin{ganttchart}%
		[canvas/.style={fill=none, draw=black!84, line width=.35pt},
		hgrid style/.style={draw=black!84, line width=.75pt},
		vgrid={{*3{draw=black!32, line width=.75pt},*1{draw=black!84, line width=.75pt}}},
		title/.style={draw=none, fill=none},
		title label font=\footnotesize,
		title label anchor/.style={below=-1.5ex},
		include title in canvas=false,
		bar label font=\footnotesize\color{black!84},
		bar label anchor/.style={left=0cm},
		bar height=0.3,
		bar/.style={draw=none, fill=black!63},
		bar incomplete/.style={fill=barblue},
		milestone progress label font=\footnotesize\color{black!70},
		group incomplete/.style={fill=groupblue},
		group left shift=0,
		group right shift=0,
		group height=.3,
		group label font=\bfseries\footnotesize,
		group label anchor/.style={left=0cm},
		link/.style={-latex, line width=1.pt, linkred},
		link label font=\footnotesize\bfseries\color{linkred}\MakeUppercase,
		link label anchor/.style={below left=-1.0ex and 0},
		x unit=0.62cm, y unit title=.35cm, y unit chart=0.34cm
		]{1}{8}
		\gantttitle{Year 1}{4}
		\gantttitle{Year 2}{4}\\
		\ganttgroup[]%
		{\textbf{Activity 1}}{1}{5} \\
		\ganttbar[]{Discovery infra.}{1}{3} \\
		\ganttbar[]{Campaigns}{2}{5} \\
		[grid]
		\ganttgroup[]%
		{\textbf{Activity 2}}{3}{7} \\
		\ganttbar[]{Hardening}{3}{6} \\
		\ganttbar[]{SBOM/secure-dev}{5}{7} \\
		[grid]
		\ganttgroup[]%
		{\textbf{Activity 3}}{5}{8} \\
		\ganttbar[]{Detection/CI}{5}{7} \\
		\ganttbar[]{Recovery+upstream}{7}{8}
	\end{ganttchart}
	\vspace{-0.4cm}
	\caption{Two-year project timeline (Track 3: up to 2 years).}
	\label{fig:timeline}
	\vspace{-2pt}
\end{wrapfigure}
Figure~\ref{fig:timeline} provides a Gantt chart for the activities outlined in this proposal.
% TODO(pesose): describe the timeline narrative and list the concrete yearly milestones the
% solicitation asks for (Track 3 "Development Plan ... including key milestones for each year").

\section{Results from Prior NSF Support}

% TODO(pesose): update each entry (award #, amount, period, IM/BI, one outcome with citation).
PI Zhao (NEU): \emph{``CAREER: SaTC: Rethinking Trusted Execution Environments for Embedded and IoT Systems"} (CNS-2237238, \$564,748, 2/2023 - 1/2028).
\ul{Intellectual Merit:} A systematic research approach to increase the trustworthiness and deployability of TEEs for embedded and IoT devices.
\ul{Broader Impacts:} Training the next generation of cybersecurity researchers, educators, and practitioners with deep theoretical understandings and practical skills.
\ul{Outcome:} Return-to-non-secure attacks on ARM Cortex-M TrustZone~\cite{ma2023dac}.

\newpage
\bibliographystyle{ieeetr}
\bibliography{bib, MyPub}
\end{document}
